\documentclass[11pt]{article}

\usepackage[a4paper,margin=2.6cm]{geometry}
\usepackage[T1]{fontenc}
\usepackage[utf8]{inputenc}
\usepackage{lmodern}
\usepackage{microtype}
\usepackage{amsmath,amssymb}
\usepackage{booktabs}
\usepackage[numbers]{natbib}
\usepackage[hidelinks]{hyperref}

% Allow line breaks inside long URLs (bibliography entries contain
% unbreakable content identifiers and query strings).
\expandafter\def\expandafter\UrlBreaks\expandafter{\UrlBreaks%
  \do\a\do\b\do\c\do\d\do\e\do\f\do\g\do\h\do\i\do\j\do\k\do\l\do\m
  \do\n\do\o\do\p\do\q\do\r\do\s\do\t\do\u\do\v\do\w\do\x\do\y\do\z
  \do\A\do\B\do\C\do\D\do\E\do\F\do\G\do\H\do\I\do\J\do\K\do\L\do\M
  \do\N\do\O\do\P\do\Q\do\R\do\S\do\T\do\U\do\V\do\W\do\X\do\Y\do\Z
  \do\0\do\1\do\2\do\3\do\4\do\5\do\6\do\7\do\8\do\9\do\=\do\?\do\&}

% Ragged-right bibliography: long URLs justify poorly.
\usepackage{etoolbox}
\apptocmd{\thebibliography}{\raggedright}{}{}

\newcommand{\tansu}{Tansu}
\newcommand{\mainnetcontract}{CDXINK2T3P46M4LWK35FVIXXHJ2XHAS4FOVCGVPJ63YV5OVTM24IY5BI}

\title{\tansu{}: decentralized project governance\\for open-source software on Stellar}

\author{%
  Pamphile T. Roy\\[2pt]
  \small Consulting Manao GmbH, Austria\\
  \small \href{https://orcid.org/0000-0001-9816-1416}{ORCID: 0000-0001-9816-1416}
}

\date{July 2026}

\begin{document}

\maketitle

\begin{abstract}
Open-source projects still lack an auditable, community-weighted record of
which commits are \emph{released} or \emph{revoked}.
CVE-2024-3094 showed maintainer capture can beat pipeline
attestation~\citep{cve2024xz,wolves2025xz}; TUF, in-toto, Sigstore, and SLSA
trace artefacts to source revisions~\citep{tuf,in_toto,newman2022sigstore,slsa}
but record no equivalent governance outcome.
\tansu{} closes that gap with Soroban contracts on Stellar
that store maintainer lists, commit hashes, weighted outcomes, and IPFS pointers
to SBOMs, CVE reports, and attestations keyed by commit. Ballots may be public
or hidden behind Pedersen commitments over BLS12-381, with the weighted tally
verified on-chain. Members may bind a Git SSH key to their Stellar account at
registration. On Stellar mainnet the contracts administer the Public Goods
Award; anonymous mode still depends on a designated tallying party for ballot
secrecy.

\medskip
\noindent\textbf{Keywords:} decentralized governance; open-source software;
software supply chain; anonymous voting; smart contracts; Stellar
\end{abstract}

\section{Introduction}
\label{sec:intro}

Most software still pulls code from a central forge. Compromise the
maintainer account or CI job and you can ship malicious releases without
touching the provenance tooling downstream users trust. Ohm et al.\ counted
174 malicious packages on npm, PyPI, and RubyGems over four
years~\citep{ohm2020backstabber}; Ladisa et al.\ later listed 107 attack
vectors across the open-source supply chain~\citep{ladisa2023taxonomy}. The
2024 XZ Utils incident is a sharper example: the attacker spent years earning
maintainer access and only then planted the
backdoor~\citep{cve2024xz,wolves2025xz}.

Artifact signing and pipeline attestation are necessary, but the XZ attack
exploited maintainer trust rather than missing signatures. Section~\ref{sec:related}
surveys TUF, in-toto, Sigstore, and SLSA; none stores a project vote that
marks a commit as \emph{released} or \emph{revoked}.

A separate risk is operational. Sigstore entries are published through operated
services such as Rekor and Fulcio~\citep{newman2022sigstore,sigstore_rekor};
TUF deployments rely on maintained metadata mirrors; the global CVE namespace
is curated under a U.S.\ government contract with
MITRE~\citep{cveprogram}. In April 2025 MITRE warned that CVE program funding
was about to lapse, which would have halted new assignments until CISA granted
a short extension~\citep{mitre_cve_funding2025}. Neither case is on-chain
fraud, but both show how much of the stack can pause when a single operator or
sponsor falters. \tansu{} does not replace those registries; it records project
votes and evidence pointers on a public ledger where, once written, the anchors
are replicated with the chain rather than held behind one vendor's uptime.

\tansu{} fills that gap with Soroban contracts written in
Rust~\citep{stellar_soroban}. Source code remains on conventional git forges
(GitHub, GitLab, Codeberg, Gitea, Bitbucket, Radicle, and others); the
ledger stores registrations, maintainer lists, commit hashes, proposals, votes,
and outcomes that any observer can read from public ledger data. We describe
contract layout, anonymous voting, commit-anchored evidence, and the Public
Goods Award deployment on mainnet.

Within those contracts, code finality attaches trust states to commits through
weighted votes; maintainers record SBOM, CVE, and attestation CIDs per commit
hash; anonymous voting hides individual choices behind Pedersen commitments
while verifying the aggregate tally on-chain; and member registration checks an
SSH signature against SEP-53. The contracts are deployed on Stellar
mainnet~\citep{stellar_expert_tansu}.

\section{Related work}
\label{sec:related}

\paragraph{Supply-chain provenance.}
The Update Framework introduced survivable key compromise for software update
systems~\citep{tuf} and underpins several package managers. in-toto extends
provenance to whole pipelines by having each step sign metadata about its
inputs and outputs~\citep{in_toto}. Sigstore removes long-lived signing keys
by binding ephemeral certificates to OpenID
identities~\citep{newman2022sigstore}, and SLSA defines provenance levels that
build systems can claim and verify~\citep{slsa}. In practice, consumers still
reach much of this material through operated registries: Sigstore's transparency
log~\citep{sigstore_rekor}, TUF metadata hosts, and the CVE program run by
MITRE under federal contract~\citep{cveprogram,mitre_cve_funding2025}. When
those services pause or their funding is uncertain, downstream tools lose a
shared reference even though the underlying signatures may still exist
offline.

\paragraph{Cryptographic voting.}
Helios demonstrated practical open-audit elections on the
web~\citep{adida2008helios}. The Open Vote Network implemented a
self-tallying, maximally private boardroom protocol as an Ethereum
contract~\citep{mccorry2017smart}. Both target a different deployment model
than a live governance contract with badge-weighted votes and collateral
deposits.

\paragraph{On-chain governance.}
Snapshot provides off-chain, gasless voting for token
communities~\citep{snapshot}, and OpenZeppelin's Governor contracts are widely
used for on-chain DAO governance on
Ethereum~\citep{openzeppelin_governor}. Fritsch et al.\ found voting power in
major DAOs concentrated in a handful of
addresses~\citep{fritsch2024dao}. Spychiger et al.\ ran a controlled DAO
experiment for project management~\citep{spychiger2023dapo}.

\tansu{} sits between the last two groups. Governance runs on-chain, and a
passed proposal can write a commit hash into contract state. Weights come from
maintainer badges or Neural Quorum Governance scores rather than token
balances. Evidence and votes are indexed by commit alongside the provenance
artefacts that TUF, in-toto, Sigstore, and SLSA help produce. Table~\ref{tab:comparison}
summarises how \tansu{} differs from adjacent governance and provenance systems.

\begin{table}[t]
\centering
\caption{Comparison with adjacent provenance and governance systems. Provenance
frameworks trace build artefacts; governance platforms record votes. Sigstore,
TUF mirrors, and the CVE program depend on operated registries whose uptime is
not guaranteed~\citep{sigstore_rekor,mitre_cve_funding2025}. \tansu{}
combines on-chain governance with commit-keyed evidence anchors.}
\label{tab:comparison}
\small
\begin{tabular}{@{}lcccc@{}}
\toprule
System & On-chain & Commit & Evidence & Private \\
       & governance & anchor & per commit & ballots \\
\midrule
TUF~\citep{tuf} & -- & -- & -- & -- \\
in-toto~\citep{in_toto} & -- & partial & attestations & -- \\
Sigstore~\citep{newman2022sigstore} & -- & partial & signatures & -- \\
Snapshot~\citep{snapshot} & off-chain & -- & -- & -- \\
OpenZeppelin Governor~\citep{openzeppelin_governor} & yes & -- & -- & -- \\
\textbf{\tansu{}} & yes & yes & yes & optional \\
\bottomrule
\end{tabular}
\end{table}

\section{System overview}
\label{sec:overview}

\tansu{} consists of a main governance contract and an auxiliary membership
contract, both written in Rust against the Soroban SDK, plus a web client,
command-line tooling, and continuous-integration recipes. The code is
released under the BSD 3-Clause License~\citep{tansu_repo}.

\noindent\textbf{Architecture.}
The forge stays the forge: Git history, pull requests, and review threads are
not copied to the chain. On-chain state is limited
to the facts a third party needs to audit governance---maintainer addresses,
the current commit hash, proposal metadata, vote weights, and results.
Proposal bodies, member profiles, and evidence files reside on IPFS while the
contract stores only their CIDs. Each state transition emits an event, so
indexers can reconstruct the full history even when the contract retains only
the latest value, such as the current commit hash.

\paragraph{Projects and commits.}
A project registers under a short name; the contract derives its key as the
Keccak-256 hash of the name and stores the maintainer set, the forge URL, and
the CID of a metadata file. Registration requires a collateral deposit, which
deters name squatting. Any maintainer can then update the project's current
commit hash. The contract keeps only the latest hash and emits an event per
update, so the full sequence of attested hashes remains recoverable from the
event stream.

\paragraph{Membership and voting weight.}
Members register themselves with a profile CID. Maintainers assign
per-project badges that determine the maximum voting weight of a member
(Table~\ref{tab:badges}). Alternatively, a proposal can be token-based, in
which case weight equals the amount of a designated token the voter locks
for the duration of the vote. A third mode delegates weights to an external
contract: for the Stellar Community Fund deployment, weights come from Neural
Quorum Governance (NQG) scores queried live from the fund's scoring contract,
with a minimum threshold that restricts voting to community
pilots~\citep{pg_award_proposer}. A separate soulbound NFT contract (SEP-50
\citep{sep50}) exposes each member's role and NQG score as token traits.

\begin{table}[t]
\centering
\caption{Badge types and their voting weights in the main contract.}
\label{tab:badges}
\begin{tabular}{lr}
\toprule
Badge & Weight \\
\midrule
Developer & $10{,}000{,}000$ \\
Triage    & $5{,}000{,}000$ \\
Community & $1{,}000{,}000$ \\
Verified  & $500{,}000$ \\
Default (no badge) & $1$ \\
\bottomrule
\end{tabular}
\end{table}

\noindent\textbf{Proposals and code finality.}
A proposal carries a title, the CID of its full text, a voting deadline, and
optionally a list of \emph{outcome contracts}: cross-contract calls executed
when the proposal is approved, rejected, or cancelled. Code finality uses this
hook. A project defines trust states for its releases (for example
\emph{released}, \emph{security-reviewed}, or \emph{revoked}) and records each
state change against a commit hash through a proposal whose outcome contract
writes the new state. The decision rule is a supermajority. With weighted
tallies $T_a$, $T_r$, and $T_{ab}$ for approve, reject, and abstain, a
proposal passes only if $T_a > T_r + T_{ab}$, is rejected only if
$T_r > T_a + T_{ab}$, and is cancelled otherwise.
Execution waits for a timelock after the voting deadline. Proposal creation and
voting require collateral deposits refunded on execution; maintainers may revoke
a malicious proposal or remove an abusive vote and withhold the collateral.
They may also bar addresses with a declared conflict of interest from voting on
a specific proposal.

\section{Commit-anchored evidence}
\label{sec:evidence}

SBOMs, CVE write-ups, and attestation bundles still reach users through
release pages, mailing lists, and vendor portals, where they are scattered and
difficult to audit after publication. Central registries compound the problem:
if Sigstore's log or the CVE feed is unreachable, consumers lose the default
discovery path even when a project already has the document. \tansu{} assigns
each project-maintained document a fixed on-chain slot keyed by the commit it
describes. The slot stores a CID and timestamp on the ledger; anyone can read
that anchor from chain history without asking a particular operator to stay
online.

The contract exposes two functions. A maintainer calls \texttt{set\_evidence}
with a project key, a commit hash, an evidence kind, and a CID; anyone calls
\texttt{get\_evidence} to retrieve the recorded history for a given commit
and kind. Three kinds are supported: \texttt{Sbom} for software
bills of materials (e.g., CycloneDX~\citep{cyclonedx} or
SPDX~\citep{spdx} documents), \texttt{Cve} for vulnerability reports
referencing CVE identifiers~\citep{cveprogram}, and \texttt{Attestation} for
build attestations such as SLSA provenance bundles~\citep{slsa}. The
evidence itself lives on IPFS; the contract stores only the CID and a
timestamp.

Evidence is append-only. Each call adds an entry to the history for the
triple (project, commit, kind) rather than overwriting it, so successive
vulnerability re-scans of the same commit accumulate: a commit that was clean
at release time and gained a CVE report two years later carries both records.
The contract keeps the most recent ten entries per triple; older entries roll
off on-chain but remain recoverable from the emitted events, which indexers
retain in full.

\noindent
In our own repository, a continuous-integration workflow generates an SBOM on
every release tag, uploads it to IPFS, and records the CID on-chain in the
same run. A single query against the contract
and an IPFS gateway returns the release commit, its SBOM, and any posted
attestations without depending on the forge remaining available or trustworthy.

The contract records who published each entry and when; only maintainers can
write. It does not parse SBOM syntax, check CVE applicability, or walk an
attestation chain, and it does not require the commit to be the project's
current head, so historical commits remain valid targets. Consumers perform
that validation off-chain. On-chain, the contract provides an append-only,
tamper-evident log of what the project asserted about each commit: later
maintainers cannot silently rewrite an earlier entry, though they can add new
ones.

\section{Anonymous voting}
\label{sec:voting}

Public voting is the default and requires no cryptography: votes are posted in
the clear and tallied by the contract. Funding rounds and personnel votes in
the Public Goods Award required ballot secrecy that public tallies could not
provide. \tansu{}'s anonymous mode never publishes individual choices, yet the
aggregate tally is still verified on-chain.

\paragraph{Setup.}
A maintainer configures anonymous voting once per project. The contract
derives two generators $G$ and $H$ of the BLS12-381 $\mathbb{G}_1$ group by
hashing two fixed labels to the curve under distinct domain-separation tags.
No discrete-logarithm relation between $G$ and $H$ is known to any party, and
no trusted setup is required. The maintainer also registers a project-level
RSA-OAEP public key (2048-bit, SHA-256) used to encrypt ballot payloads
client-side.

\paragraph{Ballots.}
A ballot is a one-hot vector $v = (v_a, v_r, v_{ab}) \in \{0,1\}^3$ over the
choices approve, reject, abstain. For each component $i$ the voter samples a
random seed $r_i$ and computes a Pedersen commitment~\citep{pedersen}
\begin{equation}
C_i = v_i\, G + r_i\, H .
\label{eq:commitment}
\end{equation}
The voter submits the three commitments together with RSA-OAEP ciphertexts of
$v$ and $r$ under the project public key. The contract checks that each $C_i$ is a
valid point in the prime-order subgroup of $\mathbb{G}_1$, stores the
ciphertexts without processing them, and folds the weighted commitments into
running aggregates
\begin{equation}
A_i \;\leftarrow\; A_i + w\, C_i ,
\label{eq:aggregate}
\end{equation}
where $w$ is the voter's weight at the time of the vote. Commitments are
computed client-side; the contract offers the same computation as a read-only
simulation helper, which must never be submitted as a transaction.

\paragraph{Tally and verification.}
After the deadline, the tallying party (a maintainer holding the decryption
key) decrypts all ballots and computes the weighted sums
$T_i = \sum_j w_j\, v_{j,i}$ and $S_i = \sum_j w_j\, r_{j,i}$. It submits
$(T, S)$ to the contract, which accepts the tally if and only if
\begin{equation}
T_i\, G + S_i\, H \;=\; A_i \qquad \text{for } i \in \{a, r, ab\}.
\label{eq:verify}
\end{equation}
By the homomorphism of \eqref{eq:commitment}, equation \eqref{eq:verify}
holds exactly when the submitted tallies are consistent with every recorded
ballot; under the discrete-logarithm assumption the commitments are binding,
so the tallying party cannot substitute a different result. The seeds provide
information-theoretic hiding, so the commitments reveal nothing about
individual choices. The supermajority rule of Section~\ref{sec:overview} is
then applied to $T$.

\paragraph{Trust model.}
Ballot privacy is not trustless in this mode. The tallying key holder can
decrypt individual votes and can stall execution by refusing to submit
$(T,S)$; lost keys freeze a proposal. The chain nonetheless enforces that any
announced tally matches the recorded commitments, which a trusted tallier
without cryptographic checks could falsify. In practice we rotate who holds
the key, and public voting remains available when the threat model permits it.

\section{Git identity binding}
\label{sec:identity}

Votes count on Stellar addresses while patches land under Git handles. Members
bind the two during \texttt{add\_member} or \texttt{update\_member} by
submitting a forge handle, an Ed25519 public key, a challenge message, and an
SSH signature over that message per SEP-53~\citep{sep53}. The contract invokes
the native Ed25519 verifier; on success it stores the handle and key and emits
an event.

The signature proves key possession at registration time. Matching the key to
the handle is left to clients: GitHub publishes
\texttt{github.com/\textless{}user\textgreater{}.keys}, GitLab provides an
equivalent endpoint, and Radicle aliases in \texttt{tansu.toml} can be checked
against repository metadata. Anyone can compare the result. The contract does not
call a forge API.

\section{Deployment and adoption}
\label{sec:deployment}

\tansu{} runs on Stellar mainnet~\citep{stellar_expert_tansu}. The contract
reports version~2; upgrades are gated by an M-of-N admin threshold and a
timelock. Development was funded through the Stellar Community Fund's
Activation and Build tracks~\citep{scf_rec_activation, scf_rec_build}.

Its main production workload is the Public Goods Award: proposal submission,
NQG-weighted voting, and outcome execution on
\tansu{}~\citep{pg_award_overview,pg_award_tansu_mods,pg_award_proposer}.
Contract changes for that process were funded
separately~\citep{scf_rec_pg}. The project also participates in Drips
Wave~\citep{drips_wave_docs}, a recurring funding programme for open-source
contributors, and each release tag in our repository triggers the evidence
workflow of Section~\ref{sec:evidence}.

\subsection{Evaluation}
\label{sec:evaluation}

The mainnet contract address is listed on Stellar
Expert~\citep{stellar_expert_tansu}; \texttt{version()} returns~2. Contract
source, client code, and integration tests are available in the public
repository~\citep{tansu_repo}. Proposals, votes, evidence entries, and commit
updates each emit events, so the history is reconstructible from ledger
archives without privileged access.

The Public Goods Award is the primary production workload on that contract.
Table~\ref{tab:pg-award} reports round outcomes verified against intake records
in the working-group repository~\citep{pg_award_intake,pg_award_overview}.
Seventeen projects received funding; total disbursement was approximately
USD~400{,}000. Individual proposal text, vote tallies, and execution
transactions remain on-chain and in the archived intake forms cited above.

\begin{table}[t]
\centering
\caption{Public Goods Award outcomes (verified from intake records in the
working-group repository~\citep{pg_award_intake}).}
\label{tab:pg-award}
\begin{tabular}{lr}
\toprule
Metric & Value \\
\midrule
Funded projects & 17 \\
Total disbursed & ${\sim}$USD~400{,}000 \\
Governance contract & mainnet v2 \\
\bottomrule
\end{tabular}
\end{table}

\section{Limitations}
\label{sec:limitations}

Anonymous voting trusts the tallying party for ballot secrecy
(Section~\ref{sec:voting}). Commit hashes reflect whatever maintainers post; a
dishonest quorum could misstate the current head, and the contract would not
detect the discrepancy unless a client compared against the forge. Evidence CIDs
are not validated, so a maintainer can anchor arbitrary content. Bulk files
still depend on IPFS pinning; the chain records pointers, not bytes, and does
not replace the CVE program or Sigstore---it complements them when operated
registries are slow or unavailable. Badge assignment, NQG scores, and NFT
admission all sit upstream of the contracts: the chain records who granted whom
weight, not whether that weight was justified.

\section{Conclusion}
\label{sec:conclusion}

\tansu{} records on Stellar what a forge cannot: weighted votes, commit
pointers, and evidence references tied to specific revisions. The Public Goods
Award deployment funded seventeen projects for roughly USD~400{,}000, with
intake records published in the working-group
repository~\citep{pg_award_intake}. Provenance frameworks trace build
artefacts to source; \tansu{} adds a governance vote and document index on the
same commit hash. Open items include threshold decryption for anonymous
voting, more evidence kinds, and finer-grained turnout statistics.

\section*{Acknowledgements}

Funded by the Stellar Community Fund (Activation and Build tracks). Thanks to
the Public Goods Award working group for running the first production
workload on the contracts. Parts of the web client and this manuscript were
edited with generative AI tools; the smart contracts were not.

\bibliographystyle{plainnat}
\bibliography{paper}

\end{document}
